\section{Discussion and Limitations}
\label{sec:discussion-residual}

The results delimit the practical contribution. At the primary 0.556\,s
deadline, Capacity returns less improvement than the shared greedy policy;
the longer 2.221\,s budget permits additional recovery, but CheapSummary
achieves higher aggregate improvement than Capacity. These comparisons do
not establish an advantage for the structural learner. The contribution is
an execution-conditioned formulation and evaluation of coordinated recovery
beyond a common feasible response. Learned delegation and heuristic
scheduling already prioritize subsolver calls using execution
outcomes~\cite{li2021delegate,chmiela2021schedule}. Action-induced
compatibility, the warm anchor, and same-state supervision define the
specific prediction problem here. Graph message passing and clique bounds
remain established components whose inclusion provides no general
performance guarantee.

Conditional recovery opportunity and timely policy improvement differ.
All policies receive the same executed greedy prefix, so its improvement
cannot be attributed to learning. The complete caller deadline charges
preparation, prediction, executor calls, membership validation, and return.
A late result retains the original schedule under the stated metric, even
when a useful candidate existed internally. The high deadline-miss rate
therefore materially affects Capacity's primary result. Native soft search
allowances describe a different boundary. Point-value divided by expected
cost also omits runtime uncertainty and is a heuristic allocation rule;
valuable cloned responses alone cannot justify its effectiveness.

Post hoc diagnostics on a fixed training subset identify sensitivity to
request representation. Structural value summaries have much larger
magnitudes than the native-budget features, while predicted residual
fractions change little across workpoints whose executed extra values
differ. Most complete fitting histories have flat conditional allocation
regret, despite parameter updates and changing auxiliary losses. Some
scalar heads approach their warm floor; this does not support a conclusion
of network-wide gradient collapse. Unit imbalance is a plausible mechanism
for weak workpoint discrimination, but these observations neither isolate
the sole cause nor establish that normalization improves timed decisions.
Offline checkpoint regret evaluates one conditional choice and does not
optimize the complete multi-call return objective.

The bounds restrict scalar values; fractional occupancies need not be
integer schedules. Feasibility follows from checking actual memberships
against original conflicts. Ownership specifies participating agents without
establishing decentralized communication or privacy. Synthetic visibility
contacts omit orbital and operational uncertainty, and public ID partitions
are artificial. Published whole-graph solvers have different initialization
and scope contracts. Generalization beyond the observed domains and fixed
coupling remains unestablished. Failures and zero-opportunity graphs retain
their weight, fits are averaged within graphs, and the current development
results do not constitute independent confirmation.
